\documentclass[14pt, a4paper]{article}

\usepackage{listings}
\usepackage{xcolor}

% ---------- Colors ----------
\definecolor{codeblue}{RGB}{0,90,160}
\definecolor{codecyan}{RGB}{0,130,140}
\definecolor{codegreen}{RGB}{60,130,75}
\definecolor{codeorange}{RGB}{180,90,30}
\definecolor{codepurple}{RGB}{130,70,150}
\definecolor{codenumber}{RGB}{120,120,120}
\definecolor{codeframe}{RGB}{190,200,210}

% ---------- C Style ----------
\lstdefinestyle{cstyle}{
    language=C,
    basicstyle=\ttfamily\fontsize{8.5}{7.5}\selectfont\color{black}, 
    keywordstyle=\color{codeblue}\bfseries,
    commentstyle=\color{codegreen},
    stringstyle=\color{codeorange},
    morekeywords={volatile},
    emph={uint8_t,uint16_t,uint32_t,uint64_t,size_t, pid_t, pthread_t},
    emphstyle=\color{codecyan},
    emph={[2]UART_STATUS_TX_EMPTY_MASK,UART_STATUS_RX_READY_MASK,pthread_create,pthread_join,malloc},
    emphstyle={[2]\color{codepurple}},
    numbers=left,
    numberstyle=\tiny\color{codenumber},
    numbersep=7pt,
    showstringspaces=false,
    showspaces=false,
    showtabs=false,
    breaklines=true,
    keepspaces=true,
    tabsize=4,
    frame=single,
    rulecolor=\color{codeframe},
    framerule=0.5pt,
    xleftmargin=18pt,
    framexleftmargin=8pt,
    columns=fullflexible,
    aboveskip=3pt,
    belowskip=3pt
}

% ---------- ARM Assembly Language ----------
\lstdefinelanguage{ARM}{
    morekeywords={LDR,LDRB,STR,STRB,MOV,ADD,SUB,AND,CMP,BEQ,BNE},
    sensitive=true,
    morecomment=[l]{@}
}

% ---------- Assembly Style ----------
\lstdefinestyle{asmstyle}{
    language=ARM,
    basicstyle=\ttfamily\fontsize{7.5}{7.5}\selectfont\color{black},
    keywordstyle=\color{codepurple}\bfseries,
    commentstyle=\color{codegreen},
    numbers=left,
    numberstyle=\tiny\color{codenumber},
    numbersep=7pt,
    showstringspaces=false,
    showspaces=false,
    showtabs=false,
    breaklines=true,
    keepspaces=true,
    tabsize=4,
    frame=single,
    rulecolor=\color{codeframe},
    framerule=0.5pt,
    xleftmargin=18pt,
    framexleftmargin=8pt,
    columns=fullflexible,
    aboveskip=3pt,
    belowskip=3pt
}

\title{CS492 Fall 2026 Unit Exam 1 Review}
\author{Sebastian Sztolberg}

\begin{document}

\maketitle

\section{Intro to MMIO}

\subsection*{Given}

\begin{itemize}
    \item UART peripheral mapped to physical memory via MMIO
    \item 16-bit address-space
    \item Base address: \texttt{0x1000}
    \item Each register is mapped to 1 byte
    \item MMIO operations take 10 cycles
    \item All other operations take 1 cycle
\end{itemize}

\vspace{3mm}

\noindent
\textbf{Register Specification}

\begin{center}
\begin{tabular}{c|c|l}
    \textbf{Offset} & \textbf{Register} & \textbf{Description} \\
    \hline
    \texttt{+0x00} & status & Status register \\
    \texttt{+0x01} & tx\_data & Transmit to port \\
    \texttt{+0x02} & rx\_data & Receive from port
\end{tabular}
\end{center}

\vspace{3mm}

\noindent
\textbf{Status Register Bit Masks}

\begin{center}
\begin{tabular}{c|l}
    \textbf{Mask} & \textbf{Meaning} \\
    \hline
    \texttt{0x01} & tx\_data empty \\
    \texttt{0x02} & rx\_data ready
\end{tabular}
\end{center}

\vspace{5mm}

\noindent
\textbf{Bit Mask Example}

\begin{center}
\begin{tabular}{c|c}
    \textbf{Value} & \textbf{Binary} \\
    \hline
    Status register ($s$) & \texttt{01010011} \\
    tx\_data mask ($t$) & \texttt{00000001} \\
    $s \mathbin{\&} t$ & \texttt{00000001}
\end{tabular}
\end{center}

\vspace{2mm}

\noindent
The result extracts the single bit representing whether
\texttt{tx\_data} is empty and ready for transmission.

\begin{itemize}
    \item Nonzero result $\rightarrow$ \texttt{tx\_data} is empty and ready to write
    \item Zero result $\rightarrow$ \texttt{tx\_data} is not empty and not ready to write
\end{itemize}

\newpage

\subsection*{Code}

\begin{lstlisting}[style=cstyle]
int main(void) {
    // Declared volatile to prevent compiler optimizations (we always want it to read/write to memory)
    volatile uint8_t * restrict status = (volatile uint8_t *)0x1000;
    volatile uint8_t * restrict tx_data = (volatile uint8_t *)0x1001;

    char c[1] = "H";

    while(!(*status & 0x01));
    *tx_data = *(volatile uint8_t *)(c);
    
    exit(EXIT_SUCCESS);
}
\end{lstlisting}

\subsection*{Corresponding Assembly}

\begin{lstlisting}[style=asmstyle]
    LDR r0, =0x1000         @ Load base address into register r0
    LDR r1, =c              @ Load character address into register r1

while_loop:
    LDR r2, [r0]            @ Load current status into register r2
    AND r2, r2, 0x01        @ Apply TX_data empty bitmask to status
    CMP r2, #0              @ Check if TX_data is empty and ready for write
    BEQ while_Loop          @ Loop if not empty

    LDR r3, [r1]            @ Load the character into register r3
    STR r3, [r0, #1]        @ Store the character in the TX_data register
    
    RET                     @ Return from the function call
\end{lstlisting}

\subsection*{Questions}

\begin{enumerate}

\item \textbf{How many cycles does the following program body take?}

\vspace{15mm}

\item \textbf{Given the fact that the CPU can process 200k cycles a second, and a full write takes 200 cycles, how many write operations do we
need to do to achieve 20 percent CPU utilization}

\vspace{15mm}

\end{enumerate}

\newpage

\section{MMIO and Cycle-Counting}

\subsection*{Given}

Assume the following specifications:

\begin{itemize}
    \item UART peripheral mapped to physical memory (MMIO)
    \item 32-bit address-space
    \item Base address: \texttt{0x10000000}
    \item Each register allocated 4 bytes
    \item Data may only be transmitted/received \textbf{1 byte} at a time
    \item MMIO operations take 10 cycles, all other operations take 1 cycle
\end{itemize}

\vspace{3mm}

\noindent
\textbf{Register Specification}

\begin{center}
\begin{tabular}{c|c|l}
    \textbf{Offset} & \textbf{Register} & \textbf{Description} \\
    \hline
    \texttt{+0x00} & status & Status register \\
    \texttt{+0x04} & tx\_data & Transmit to port \\
    \texttt{+0x08} & rx\_data & Receive from port
\end{tabular}
\end{center}

\vspace{3mm}

\noindent
\textbf{Status Register Bit Masks}

\begin{center}
\begin{tabular}{c|l}
    \textbf{Mask} & \textbf{Meaning} \\
    \hline
    \texttt{0x01} & tx\_data empty \\
    \texttt{0x02} & rx\_data ready
\end{tabular}
\end{center}

\vspace{5mm}

\newpage

\subsection*{Code}

\begin{lstlisting}[style=cstyle]
#define UART_STATUS_TX_EMPTY_MASK 0x01
#define UART_STATUS_RX_READY_MASK 0x02

int main(void) {
    volatile uint32_t * uart_status =
        (volatile uint32_t*)(0x10000000);

    volatile uint32_t * uart_tx_data =
        (volatile uint32_t*)(0x10000000 + 0x04);

    volatile uint32_t * uart_rx_data =
        (volatile uint32_t*)(0x10000000 + 0x08);

    uint8_t w[8];

    for (int i = 0; i < 8; i++) {
        while((*uart_status & UART_STATUS_RX_READY_MASK) == 0);

        w[i] = *(uint8_t*)(uart_rx_data);
    }

    for (int j = 0; j < 8; j++) {
        while((*uart_status & UART_STATUS_TX_EMPTY_MASK) == 0);

        *uart_tx_data = *(volatile uint32_t*)(&(w[j]));
    }

    return 0;
}
\end{lstlisting}

\vspace{5mm}

\subsection*{Corresponding Function Body Assembly:}

\begin{lstlisting}[style=asmstyle]
    LDR r1, =0x10000000     @ base address
    LDR r2, =w              @ word array address

    MOV r0, #8              @ loop 1 counter

for_loop1:
while_loop1:
    LDR r3, [r1]
    AND r3, r3, 0x02
    CMP r3, #0
    BEQ while_loop1

    LDRB r3, [r1, #8]
    STRB r3, [r2]

    ADD r2, r2, #1
    SUB r0, r0, #1
    CMP r0, #0
    BNE for_loop1


    LDR r2, =w              @ reload word arr base addr
    MOV r0, #8              @ reset loop counter


for_loop2:
while_loop2:
    LDR r3, [r1]            @ load status register
    AND r3, r3, 0x01        @ apply tx_data bitmask
    CMP r3, #0               @ cmp
    BEQ while_loop2         @ keep iterating if not ready

    LDRB r3, [r2]           @ read w[i]
    STRB r3, [r1, #4]       @ store in tx_data

    ADD r2, r2, #1
    SUB r0, r0, #1
    CMP r0, #0
    BNE for_loop2
\end{lstlisting}

\subsection*{Questions}

\begin{enumerate}
    \item \textbf{What is the program doing?}

    \vspace{15mm}

    \item \textbf{Under the assumption that both while loops take the same
    amount of cycles, will both for loops take the same amount of cycles?
    Show your work.}

    \vspace{25mm}

    \item \textbf{If for each for-loop iteration, the while loop iterates
    10 times, how many cycles will it take to complete the function body?}

    \begin{center}
        \textit{Assume the successful read/write occurs on the 11th cycle.}
    \end{center}

    \vspace{30mm}

    \item \textbf{How would we implement a function that constantly waits
    to receive new data from the port, and then transmits it to a 1-byte
    software buffer that can be read by some external process?}

    Give C-style pseudocode based on the C code above.

    \vspace{35mm}
\end{enumerate}

\newpage

\section{Fork Behavior}

\begin{lstlisting}[style=cstyle]
int main(int argc, char* argv[]) {
    pid_t p1;
    pid_t p2;

    uint8_t x = 0;

    if ((p1 = fork())) {
        x++;
    }

    printf("%u\n", x);

    if (p1 && (p2 = fork())) {
        x++;
    }

    printf("%u\n", x);

    return 0;
}
\end{lstlisting}

\subsection*{Questions}

\begin{enumerate}
\item \textbf{Upon a call to fork(), what elements of the virtual memory image are shared between the parent and child? Explain.}

\vspace{15mm}

\item \textbf{How many processes are there in total? (include the parent)}

\vspace{15mm}

\item \textbf{List the prints produced by each process individually}

\vspace{15mm}

\end{enumerate}

\newpage

\section{Fork, Exec, Wait}

\begin{lstlisting}[style=cstyle]
int main(void) {
    pid_t p1;
    pid_t p2;
    
    int n = 0;

    if ((p1 = fork())) {
        n++;
    }

    printf("%d\n", n);

    if (!p1) {
        p2 = fork();
        n += 2;
    }

    if (!p2) {
        execvp("./helper", NULL);
    }

    if (p1) {
        wait(NULL);
        printf("%d\n", n);
    } else {
        wait(NULL);
        printf("%d\n", n);
    }

    exit(EXIT_SUCCESS);
}
\end{lstlisting}

\begin{lstlisting}[style=cstyle]
// helper.c code
int main(void) {
    printf("%d\n", 3);
    exit(EXIT_SUCCESS);
}
\end{lstlisting}

\subsection*{Questions}

\begin{enumerate}

\item \textbf{Explain what execvp() and wait() do}

\vspace{15mm}

\item \textbf{How many processes are spawned during the execution of this program? (include the parent)}

\vspace{15mm}

\item \textbf{Please give the independent strings printed by each process}

\vspace{15mm}

\item \textbf{Bonus: what do the suffixes of exec signify? i.e. [v, l] and [p, e]}

\vspace{15mm}

\end{enumerate}

\newpage

\section{PThreads and Thread Shared Resources}

\begin{lstlisting}[style=cstyle]
int main(void) {
    pthread_t t1;
    pthread_t t2;

    int* np = (int*)malloc(sizeof(int));
    *np = 0;


    pthread_create(&t1, NULL, helper, (void*)np);
    pthread_create(&t2, NULL, helper, (void*)np);

    pthread_join(t1, NULL);
    pthread_join(t2, NULL);

    *np += 1;

    printf("%d\n", *np);

    free(np);
    exit(EXIT_SUCCESS);
}
\end{lstlisting}

\begin{enumerate}

\item \textbf{What are some differences between processes and threads?}

\vspace{15mm}

\item \textbf{What values are printed by the program above? Assume atomicity per line of C code.}

\vspace{15mm}

\item \textbf{We would like to truly parallelize our program. 25 percent of our program is parallelizable. We have 4 CPU cores at our disposal. Please calculate the CPU speedup.}

\vspace{15mm}

\end{enumerate}

\end{document}